\section*{零基础阅读准备：先学会读一句数学话}
这份导论没有要求你立刻证明一个深定理。它解释为什么整门课要从集合出发，最后走到经济均衡。你目前最需要的是会分辨四种句子：\textbf{定义}约定词语是什么意思；\textbf{假设}规定这道题允许使用什么；\textbf{结论}是要推出什么；\textbf{证明}把前两者一步步连到结论。

记号 $x\in A$ 读作“$x$ 是集合 $A$ 的元素”；$A\subseteq B$ 读作“$A$ 的每个元素都属于 $B$”。$\forall$ 是“对每一个”，$\exists$ 是“至少存在一个”。“若 $P$，则 $Q$”写成 $P\Rightarrow Q$；它\emph{不}自动说明 $Q\Rightarrow P$。$P\Leftrightarrow Q$ 则要求两个方向都成立。

\paragraph{一个最小证明范例。}若 $A\subseteq B$ 且 $B\subseteq C$，证明 $A\subseteq C$。任取 $x\in A$；由第一个假设 $x\in B$，由第二个假设 $x\in C$；因为 $x$ 是任取的，所有 $A$ 中的元素都在 $C$ 中，故 $A\subseteq C$。这四句话已经包含本课大量集合证明的骨架。

\Slide{1}{导论的对象与参考书}
本页是课程的标题页：课件围绕 Gerard Debreu 的 \emph{Theory of Value: An Axiomatic Analysis of Economic Equilibrium}（1959），授课人为张顺明教授。这里的 value 处于一般均衡理论语境，不是投资估值课程。“公理化分析”指先列清楚对象和假设，再用逻辑推导结论。

本册对应教材第一章第 1 节 Introduction。原书正文很精练，因此本套 notes 会在原课件截图后补上中间步骤；讲解是自学辅助，不是老师额外布置的考纲。

\Slide{2}{全章地图：十个主题怎样接起来}
章节顺序是：导论、集合、函数与对应、预序、实数、$\R^N$ 中的极限、连续函数、连续对应、向量、固定点。可以读成三条依赖链：
\[
\text{集合}\longrightarrow\text{函数/对应}\longrightarrow\text{连续性},\qquad
\text{实数}\longrightarrow\text{极限}\longrightarrow\text{紧性},
\]
\[
\text{向量与凸性}+\text{紧性}+\text{连续性}\longrightarrow\text{固定点}.
\]
预序描述“哪个至少不差于哪个”，为偏好关系作准备。$\R^N$ 表示由 $N$ 个实数坐标组成的点的集合；例如两个商品的消费量 $(x_1,x_2)$ 就是 $\R^2$ 的一个点。学习时先掌握箭头左端，后面就不需要靠死记。

\Slide{3}{“原则上不需要数学基础”应怎样理解}
原书说第一章提供后文要使用的数学概念与结果，并说“原则上”不需已有数学知识。这是对内容组织的说明，不表示第一次读就应该看懂。所谓从零开始，是定义尽量自给自足；但把定义用在证明中仍需练习。

\paragraph{本页的实际学习任务。}读到一个新词时写下：（1）它描述什么对象；（2）有哪些条件；（3）一个满足条件的例子；（4）一个只差一个条件的反例。例如“非空集合”只说至少有一个元素，$\{0\}$ 非空，而 $\varnothing$ 为空；$0$ 与 $\{0\}$ 也不是同一对象。

原文括号说明书末 notes 中会自由用到其他知识；因此不要把“正文从基础开始”误读为“所有书末补充都适合第一次阅读”。

\Slide{4}{概念的逻辑顺序与先修关系}
本页逐节解释安排。集合、函数/对应、预序先保持抽象：不必一开始就假定元素都是数字。函数给每个输入一个输出；对应允许一个输入得到一整个集合。预序在集合上比较元素，不必能够比较任何两者。

实数一节是关键转折。定义实数前先讨论\textbf{二元运算}（把一对元素变成一个元素，例如加法）和\textbf{最小上界}（所有上界中最小的一个）。之后由 $\R$ 构造 $\R^N$，用点列研究“越来越靠近”，再定义连续性。向量结构额外允许加法和数乘；固定点要求输入属于自己的输出集合。

\paragraph{把术语放进一个例子。}消费者可行集合 $B$ 是集合；效用 $u:B\to\R$ 是函数；最优选择 $D(p)$ 可能不唯一，故它适合写成对应；若每个价格 $p$ 的可行集合与目标函数满足适当条件，才可能证明 $D$ 有连续性性质。这里我们只是建立词汇联系，尚未断言均衡一定存在。

\Slide{5}{公理化方法为什么能帮助经济分析}
原文强调：清楚地定义概念，能帮助正确提出并解决经济问题。证明题里最常见的错误不是算错，而是把没有给出的条件偷偷用上。例如“价格上升，所以需求必定下降”并非不加任何偏好假设就成立的数学定理。

\paragraph{可照着使用的写法。}先写对象 $S$，再写条件“$S$ 非空、紧、凸”，最后写结论。每次使用一个定理，都核对它要求的每个条件。例如连续函数在\emph{非空紧集}上能取到最大值；若只记“连续函数有最大值”，$f(x)=x$ 在 $(0,1)$ 上就给出反例。$f(x)<1$ 对所有 $x$ 成立，但总能找到更大的 $x$，所以没有最大值点。

\Slide{6}{作者承认的实际门槛：抽象思考需要练习}
本页特别承认：“不需先修”只在原则上成立。实际需要抽象思考、短时间吸收新概念的能力；还默认了基本集合逻辑与整数知识。你不必把不熟练理解成能力问题，应该把它拆成可以练的技能。

\paragraph{三个具体动作。}第一，把抽象变量换成小例子，如 $A=\{1,2\}$。第二，再恢复一般记号，检查推理是否仍适用于任意 $A$。第三，用反例测试自己漏掉的条件。例子只能帮助理解，\emph{不能代替一般证明}；验证 $1+2=2+1$ 不等于证明所有实数加法交换律。

\paragraph{量词次序练习。}“对每个实数 $x$，存在 $y>x$”成立，取 $y=x+1$；“存在一个实数 $y$ 大于所有实数 $x$”不成立，因为取 $x=y+1$ 即矛盾。两句只交换了 $\forall$ 与 $\exists$，意思已经完全不同。

\Slide{7}{数学工具与经济结论的三条对应}
原页给出三个具体例子。第一，$\R^N$ 的每个子集都有\textbf{可数稠密子集}，将用于效用函数存在性论证。稠密的直觉是：集合内任一点附近都有这个子集的点；可数是能够用自然数编号，不表示点之间有最小距离。第二，凸集的支撑超平面定理是最优性分析的重要工具：一条直线或一个高维平面可以在边界“托住”凸集。第三，Kakutani 固定点定理支持一般均衡存在性证明。

\paragraph{不要倒读成无条件结论。}“工具被用于证明均衡”不等于“任意经济都有均衡”；非空性、紧性、凸性、连续性等数学假设必须从经济模型的假设中逐一导出。考试写固定点应用时，最重要的通常是\emph{验证假设}，不是最后写一句“由 Kakutani 得证”。

教材原文采用 $1.6.n$、$1.9.x$、$1.10.d$ 这类定位：分别是第一章第 6、9、10 节的字母段落，并非乘法或小数。

\Slide{8}{这些基础工具的用途超出本书}
本页说这些概念大多属于数学中最基础的工具，应用远超本书。这里的“基础”是指后续理论依赖它们，不是指它们必然容易。集合、连续性、紧性和凸性同样用于优化、概率与博弈论。

\paragraph{这对复习的意义。}优先掌握可以迁移的证明动作：证明集合相等要证双向包含；证明存在极值常取逼近上确界的序列再抽收敛子列；证明连续常使用逆像开集或点列。把“题名”换成“证明动作”来归纳，比逐题硬背更能应对新题。

\Slide{9}{原书的跳读建议与证明考试的区别}
作者考虑到读者时间有限，把经济章节中的概念与结果写得醒目，让准备较少的读者也可理解经济含义。这是原书提供的阅读选择，而本课程以证明为考核，因此只懂直觉还不够。

\paragraph{适合你的两遍阅读法。}第一遍先看原页、中文概念解释与例子，知道定理要说什么。第二遍遮住证明，依次写出“已知、要证、选取的对象、用到的性质、结束语”。若卡住，定位到底是不懂某个定义，还是不知道怎样选择一个序列、邻域或中间量；只回查这一个环节，随后重新独立写一遍。

这一安排是学习建议，不表示任何原页或证明可以在本册中省去。

\Slide{10}{原书不提供证明：怎样把陈述变成可写的证明}
原页明确说本章通常不给断言的证明，鼓励读者重构；遇到重构较难的结论，会写“One can prove...”等提示。小字通常是例子或启发性说明，不承担正文逻辑推导。

\paragraph{你的证明清单。}看到定理后，先逐字抄出假设和结论。要证“所有”，从任意元素开始；要证“存在”，构造对象或调用一个已验证条件的存在性定理；要证“当且仅当”，分别证明两个方向；要证不可能，先假设相反结论并推到矛盾。最后检查结论是否覆盖了原命题中的全部量词。

\paragraph{一份完整的小证明。}命题：若 $A\subseteq B$，则 $A\cap C\subseteq B\cap C$。任取 $x\in A\cap C$，由交集定义有 $x\in A$ 且 $x\in C$；再由假设得 $x\in B$；于是 $x\in B\cap C$。任意性给出所求包含。这份证明没有数字计算，却是数学证明的标准形式。

\section*{本册自测与答案}
\begin{enumerate}
\item \textbf{题：}“若 $x>2$ 则 $x>0$”的逆命题是什么，是否成立？\textbf{答：}逆命题是“若 $x>0$ 则 $x>2$”，取 $x=1$ 即否定。原命题与逆命题必须分开判断。
\item \textbf{题：}为什么一个例子不能证明“所有连续函数都有最大值”？\textbf{答：}“所有”覆盖每个对象，单个例子只验证一个对象；而一个反例足以否定全称命题。例如 $f(x)=x$ 在 $\R$ 上连续且无最大值。
\item \textbf{题：}证明 $A=B$ 应写几部分？\textbf{答：}证明 $A\subseteq B$ 与 $B\subseteq A$ 两部分；每部分都可从任意元素出发。只写一边通常不够。
\item \textbf{题：}第 7 页的三条经济应用意味着所有经济都存在均衡吗？\textbf{答：}不意味着。必须先由具体经济假设验证相应数学定理的全部条件；工具的存在与工具适用于某个模型是两回事。
\end{enumerate}
